\section{Partial-Weyl recognition beyond AME states}
\label{subsec:recognition}

The full-Weyl criterion recovers the entire local Weyl basis from one
marginal.  More generally, several marginals may each recover only part of
that basis, while their labels together generate it.  The following
criterion isolates this mechanism.  It uses only a finite abelian label
group and a Weyl basis whose multiplication is additive up to phase.

Call an operator on \(r\) parties \emph{partial-Weyl diagonal} if there are
a finite index set \(P\) with a distinguished element \(0\), injections
\(f_i:P\to\F_q^2\) with \(f_i(0)=0\), and nonzero coefficients, for which
\[
 R=\sum_{v\in P}\lambda_v\bigotimes_{i=1}^{r}W_{f_i(v)},
 \qquad \lambda_v\ne0.
\]
Write \(N=|P|\).  The full-Weyl condition \((3.2)\) is the case
\(N=q^2\), where the injections are forced to be bijections.

\begin{proposition}[Partial-Weyl marginal criterion]
\label{prop:partial-weyl-marginal}
Let \(R\) and \(R'\) be partial-Weyl diagonal operators on the same
\(r\) parties with index sets of the same size \(N\geq2\), possibly with
different index sets, injections, and coefficients, and suppose
\[
 R'=(U_1\otimes\cdots\otimes U_r)R(U_1\otimes\cdots\otimes U_r)^\dagger .
\]
Assume \(r\geq3\), or \(N=2\).  Then for each \(i\), conjugation by
\(U_i\) carries every \(W_{f_i(v)}\) to a phase times a Weyl operator, and
induces a bijection \(f_i(P)\to f_i'(P')\) fixing \(0\).
\end{proposition}

\begin{proof}
Suppose \(r\geq3\).  Let \(E_i\) be the span of
\(\{W_{f_i(v)}:v\in P\}\) inside the single-party Hilbert--Schmidt space
and \(E_i'\) the span of \(\{W_{f_i'(v')}:v'\in P'\}\).  Both have
dimension \(N\), because the Weyl operators are linearly independent and
the injections are injective, and the rescaled operators
\(q^{-1/2}W_v\) are orthonormal.  All \(N\) coefficients of \(R\) are
nonzero, so by the rank computation in the proof of
Lemma~\ref{lem:diagonal-axes} the smallest subspace of factor \(i\)
supporting \(R\) is \(E_i\), and likewise for \(R'\) and \(E_i'\); hence
the product of the conjugations carries \(E_i\) onto \(E_i'\).  Now apply
Lemma~\ref{lem:diagonal-axes} to the two tensors in
\(E_1\otimes\cdots\otimes E_r\) and \(E_1'\otimes\cdots\otimes E_r'\).
Its proof uses only \(r\geq3\), nonzero coefficients, and orthonormal
bases, never \(N=q^2\).  The maps are therefore monomial in the displayed
bases, and each scalar is a phase because conjugation is unitary for the
Hilbert--Schmidt inner product.  The identity axis is fixed, so the
bijection fixes \(0\).

Suppose instead \(N=2\), with \(R=\lambda_0I+\lambda_1\bigotimes_iW_{f_i(1)}\)
and \(R'=\lambda_0'I+\lambda_1'\bigotimes_iW_{f_i'(1)}\).  Conjugation
fixes the identity, so comparing traces gives \(\lambda_0=\lambda_0'\);
subtracting leaves an equality of nonzero product operators, whose tensor
factorization gives
\(U_iW_{f_i(1)}U_i^\dagger=c_iW_{f_i'(1)}\) with \(\prod_ic_i=\lambda_1'/\lambda_1\),
and each \(c_i\) is a phase by unitarity.
\end{proof}

The arity condition cannot be removed in this form.  For two parties and
three or more terms, equal coefficient moduli leave a unitary freedom among
the corresponding Schmidt axes.  With two terms, conjugation fixes the
identity term and hence singles out the other axis.

Fix a product unitary carrying \(\lvert\psi\rangle\) to \(\lvert\phi\rangle\)
up to phase, after absorbing any party permutation by relabelling
\(\lvert\phi\rangle\).  For each party \(i\) put
\[
 \Lambda_i=\{a\in\F_q^2:\ U_iW_aU_i^\dagger\in\mathbb C^\times W_b
   \text{ for some }b\},
\]
the \emph{recognition group} at \(i\).

\begin{lemma}[Recognition group]\label{lem:recognition-group}
\(\Lambda_i\) is a subgroup of \(\F_q^2\), the induced map \(a\mapsto b\)
is an injective additive homomorphism, and \(U_i\) is Clifford if and only
if \(\Lambda_i=\F_q^2\).
\end{lemma}

\begin{proof}
If \(U_iW_aU_i^\dagger=\alpha W_{a'}\) and \(U_iW_bU_i^\dagger=\beta W_{b'}\),
then \(W_aW_b\) is a phase times \(W_{a+b}\), so conjugating that product
shows \(U_iW_{a+b}U_i^\dagger\) is a scalar times \(W_{a'+b'}\).  Hence
\(\Lambda_i\) is closed under addition and the induced map is additive; it
is injective because distinct labels have distinct projective axes.  If
\(\Lambda_i=\F_q^2\) the induced map is an injective endomorphism of a
finite group, hence bijective, so conjugation permutes all projective Weyl
axes, which is the Clifford condition.  The converse is immediate.
\end{proof}

When \(q=2\), a recognition group containing a nonzero label is either a
line or all of \(\F_2^2\).  The criterion below assembles labels recovered
from different marginals; in larger local dimension it also accommodates
intermediate subgroups that have no qubit analogue.

\begin{corollary}[Generation criterion]\label{cor:recognition-generation}
Suppose that for each party \(i\) there is a family of party sets
containing \(i\), each with at least three parties or with a two-element
index set, such that \(\rho^\psi_S\) and \(\rho^\phi_S\) are partial-Weyl
diagonal with index sets of equal size.  If the labels recovered at \(i\)
over that family generate \(\F_q^2\), then \(U_i\) is Clifford.
\end{corollary}

\begin{proof}
Partial trace gives a product-unitary equivalence of the two reduced
operators on each set.  Proposition~\ref{prop:partial-weyl-marginal} places
each recovered label in \(\Lambda_i\), and
Lemma~\ref{lem:recognition-group} closes \(\Lambda_i\) under addition and
converts \(\Lambda_i=\F_q^2\) into the Clifford property.
\end{proof}

Corollary~\ref{cor:full-weyl-cover} is the case of one set per party with
\(N=q^2\), where generation is automatic.  For stabilizer states the
supply of usable sets is exactly the minimal supports.

\begin{lemma}[Minimal supports carry partial-Weyl marginals]
\label{lem:minimal-support-charts}
Let \(\lvert\psi\rangle\) be a stabilizer state and \(\omega\) a minimal
support of its stabilizer, and let \(L(\omega)\) be the group of labels
supported inside \(\omega\).  Then for every \(i\in\omega\) the coordinate
projection \(\pi_i:L(\omega)\to\F_q^2\) is injective.  Hence \(L(\omega)\)
is isomorphic to a subgroup of \(\F_q^2\), every nonzero element of
\(L(\omega)\) has support exactly \(\omega\), and \(\rho_\omega\) is
partial-Weyl diagonal with index set \(L(\omega)\) and injections
\(\pi_i\).  Conversely, if the retained label group on a party set has at
least two elements and injects at every party, that set is a minimal
support.
\end{lemma}

\begin{proof}
An element of the kernel of \(\pi_i\) has support contained in
\(\omega\setminus\{i\}\), a proper subset of \(\omega\), so it is zero by
minimality.  The partial trace of \(\lvert\psi\rangle\langle\psi\rvert\)
over the complement of \(\omega\) retains exactly the stabilizer elements
supported inside \(\omega\), each with a coefficient of modulus one.  For
the converse, injectivity at every party of the set says that no nonzero
element of its label group vanishes anywhere on it, so no nonzero element
has smaller support.
\end{proof}

At \(q=2\) the possible orders of \(L(\omega)\) are \(2\) and \(4\), which
is the dichotomy of \cite[Lemma~1]{VanDenNestDehaeneDeMoor2005}; the
intermediate orders available for \(q>2\) are invisible there.  They occur
naturally in CSS states: a minimum-weight \(X\)-support carrying no
\(Z\)-word contributes a line of \(q\) labels.

For two LU-equivalent stabilizer states, marginal purity preserves both the
minimal supports and the order of each retained label group, since
\[
 \operatorname{Tr}(\rho_\omega^2)
   =\frac{|L(\omega)|}{q^{|\omega|}}.
\]
Thus every source chart above has a target chart of the same size.  In the
qubit case, the hypothesis of
\cite[Theorem~1]{VanDenNestDehaeneDeMoor2005} says exactly that the labels
from all minimal supports through a party generate \(\F_2^2\).  Label
groups of order two use the two-term case of
Proposition~\ref{prop:partial-weyl-marginal}; label groups of order four
use its \(r\geq3\) case.  Their standing assumptions of at least three
parties and full entanglement exclude a two-party support in the latter case.
Corollary~\ref{cor:recognition-generation} therefore recovers their full
theorem, not only its single-support corollary.

\begin{corollary}[CSS states over a prime field]\label{cor:css-recognition}
Let \(q\) be prime and let \(\lvert\psi\rangle\) be a CSS coset state such
that every coordinate lies in two minimal supports of the stabilizer, each
with at least three coordinates, one whose label group contains a nonzero
\(X\)-type label and one whose label group contains a nonzero \(Z\)-type
label.  Then every product-unitary equivalence to another such state has
Clifford factors.
\end{corollary}

\begin{proof}
The preceding purity argument supplies matching partial-Weyl marginals for
the target state.
The two supports contribute independent labels \((x,0)\) and \((0,z)\) to
\(\Lambda_i\), which generate \(\F_q^2\) over a prime field.  Apply
Corollary~\ref{cor:recognition-generation}.
\end{proof}

A minimum-weight codeword of the \(X\) code need not give a minimal support
of the stabilizer, since a \(Z\)-word may sit strictly inside it; the
hypothesis above is about minimal supports of the full stabilizer.  What is
not automatic for an arbitrary additive stabilizer is either coverage by
minimal supports or the availability of both label types at every party;
the corollary keeps both requirements explicit.

\begin{remark}[Arbitrary local dimension]\label{rem:arbitrary-dimension}
Nothing in the partial-Weyl condition,
Proposition~\ref{prop:partial-weyl-marginal},
Lemma~\ref{lem:recognition-group},
Corollary~\ref{cor:recognition-generation}, or
Lemma~\ref{lem:minimal-support-charts} uses the field structure of
\(\F_q\).  Each uses only that the labels form a finite abelian group whose
Weyl operators are an orthogonal Hilbert--Schmidt basis multiplying
additively up to phase.  Replacing \(\F_q^2\) by \(\mathbb Z_d^2\), with
\(X\) the shift and \(Z\) the clock operator on \(\mathbb C^d\), all five
statements hold verbatim at every local dimension \(d\), prime power or
not.  At composite \(d\), a unitary whose conjugation permutes the
projective Weyl axes still normalizes the Pauli group itself and is
Clifford in the standard sense.  Write \(UW_aU^\dagger=cW_{a'}\).  For
odd \(d\), every Weyl operator satisfies \(W_a^d=I\); taking \(d\)-th powers
gives \(c^d=1\).  For even \(d\), every Weyl operator satisfies
\(W_a^{2d}=I\), and the standard Pauli convention includes the \(2d\)-th
roots of unity; taking \(2d\)-th powers gives \(c^{2d}=1\).  Thus \(c\) is
already a Pauli phase and \(UW_aU^\dagger\) lies in the Pauli group.  When
\(d\) is a proper prime power, the
\(\mathbb Z_d\) Weyl system and the \(\F_d\) Weyl system on \(\mathbb C^d\)
are different, and a criterion instantiated in one says nothing in the
other.
\end{remark}

The generation hypothesis is sufficient, not necessary.  Wong and Jiang
exhibit at \(d=9\) a three-party stabilizer state that is LU but not LC
equivalent to the GHZ state \cite{WongJiang2025}.  Its three two-party
minimal supports each carry nine labels whose projections lie in
\(3\mathbb Z_9\times3\mathbb Z_9\).  They generate a subgroup of index
nine rather than \(\mathbb Z_9^2\), so
Corollary~\ref{cor:recognition-generation} does not apply.

Chang and Jing introduced the general Kruskal-type uniqueness method for
coefficient tensors in a discrete operator basis into the qubit LU problem
\cite{ChangJing2022}.  Their generic Pauli tensors are not diagonal in that
basis, and uniqueness yields invariants separating LU classes rather than a
Clifford conclusion.  For qubit stabilizer states that conclusion is
\cite[Theorem~1]{VanDenNestDehaeneDeMoor2005}.
